% =====================================================================
%  CV COMPLET — Nicolás Plata Molano · format une colonne
%  (version détaillée ; la version courte d'une page est CV_Nicolas_Plata_Molano_FR.tex)
%  Compilation : pdflatex CV_Nicolas_Plata_Molano_FR_complet.tex (deux passes)
% =====================================================================
\documentclass[10pt,a4paper]{article}

\usepackage[utf8]{inputenc}
\usepackage[T1]{fontenc}
\usepackage[french]{babel}
\usepackage[sfdefault,book,lining]{FiraSans}
\usepackage[left=1.6cm,right=1.6cm,top=1.3cm,bottom=1.6cm]{geometry}
\usepackage{xcolor}
\usepackage{titlesec}
\usepackage{enumitem}
\usepackage{tabularx}
\usepackage[hidelinks]{hyperref}
\usepackage[expansion=false]{microtype}
\usepackage{amssymb}
\usepackage{fancyhdr}
\usepackage{lastpage}

\definecolor{navy}{HTML}{1F3F6E}
\definecolor{grey}{HTML}{5C6770}

\pagestyle{fancy}
\fancyhf{}
\renewcommand{\headrulewidth}{0pt}
\fancyfoot[L]{\scriptsize\color{grey}Nicolás Plata Molano — CV détaillé}
\fancyfoot[R]{\scriptsize\color{grey}\thepage\,/\,\pageref{LastPage}}

\setlength{\parindent}{0pt}
\setlength{\parskip}{0pt}
\linespread{1.05}

\titleformat{\section}{\color{navy}\large\bfseries}{}{0pt}{}[\vspace{1pt}{\color{navy}\titlerule[0.6pt]}]
\titlespacing*{\section}{0pt}{11pt}{5pt}
\titleformat{\subsection}{\color{navy}\bfseries}{}{0pt}{}
\titlespacing*{\subsection}{0pt}{6pt}{3pt}

% \entree{Titre — rôle}{Lieu / dates}{texte}
\newcommand{\entree}[3]{%
  \textbf{#1}\hfill{\small\itshape #2}\par\vspace{1pt}
  {\small #3}\par\vspace{6pt}}
\newenvironment{points}{\begin{itemize}[leftmargin=1.1em,itemsep=1pt,topsep=2pt,parsep=0pt,label=\textcolor{navy}{\scriptsize$\blacksquare$}]\small}{\end{itemize}}
\newcommand{\outils}[1]{{\small\color{grey}\textit{Outils :} #1}}
\newcommand{\bati}{BaTiO\textsubscript{3}}

\begin{document}
\thispagestyle{fancy}

% ======================= EN-TÊTE =======================
\begin{center}
  {\LARGE\bfseries NICOLÁS PLATA MOLANO}\\[2pt]
  {\color{grey}Élève-ingénieur en double diplôme — IMT Atlantique (Automatique \& Systèmes Cyber-physiques) · Universidad Nacional de Colombia (Mécatronique)}\\[3pt]
  {\small \href{mailto:nicao511ng@gmail.com}{nicao511ng@gmail.com} \;|\; +33 7 45 57 71 66 \;|\; Nantes, France \;|\; Nationalité colombienne, titre de séjour étudiant}\\[1pt]
  {\small \href{https://www.linkedin.com/in/nicolas-plata-molano}{linkedin.com/in/nicolas-plata-molano} \;|\;
  \href{https://github.com/NicolasPlata2004}{github.com/NicolasPlata2004} \;|\;
  \href{https://nicolasplata2004.github.io}{\textbf{Portfolio : nicolasplata2004.github.io}}}
\end{center}
\vspace{-4pt}
{\color{navy}\rule{\linewidth}{0.8pt}}
\vspace{3pt}

{\small\textbf{Objectif — stage de 2\textsuperscript{e} année} : 4 à 5 mois à partir d'avril 2027 (ou 2 à 3 mois dès juin) en automatique,
robotique, systèmes embarqués ou automatisation industrielle. Nantes en priorité, mobile dans toute la France.}

% ======================= PROFIL =======================
\section{Profil}
{\small Élève-ingénieur en mécatronique, \textbf{5\textsuperscript{e} de ma promotion} (moyenne 4,4/5), admis en double diplôme à l'IMT Atlantique.
Mon fil conducteur : passer du modèle au système réel. J'identifie des procédés, je conçois leurs lois de commande, je les implémente sur
microcontrôleur ou sur automate, puis je les valide sur banc : quatre lois de commande sur un aéropendule, un bras robotique 2R qui dessine,
un objet IoT connecté en MQTT et une machine automatisée validée par jumeau numérique. En recherche, je contribue à des capteurs
piézoélectriques sans plomb (\bati) au sein du groupe KYMA. J'aime aussi expliquer : je réalise des animations pédagogiques (Manim)
et j'ai fondé un blog de vulgarisation scientifique.}

% ======================= FORMATION =======================
\section{Formation}
\entree{IMT Atlantique — Diplôme d'ingénieur généraliste (FISE 2A), double diplôme}{Nantes, 2026–2028}{%
Thématique d'approfondissement \textbf{ASCY — Automatique et Systèmes Cyber-physiques}. Diplôme accrédité par la CTI.
Enseignements : asservissements monovariables robustes, représentation d'état.
Commande d'entreprise en cours avec la startup NEXERGY (voir « Projets »).}
\entree{Universidad Nacional de Colombia — Ingénierie mécatronique}{Bogotá, 2021–2026}{%
\textit{Moyenne :} \textbf{4,4 / 5} \;|\; \textbf{5\textsuperscript{e} de la promotion} \;|\; \textit{Cursus validé :} 73,7 \% \;|\; Moniteur académique.\\
Cours marquants : Techniques de commande (PID, H$\infty$, MRAC, mode glissant), Servomécanismes, Projet industriel, Microcontrôleurs,
Électronique numérique, Réseaux informatiques, Instrumentation industrielle, Robotique, Procédés de fabrication II (4,8/5), Structures de données.}
\entree{I.E. Camilo Torres Restrepo — Baccalauréat}{Aguazul, Casanare}{%
Meilleure moyenne de l'établissement ; score ICFES remarqué au niveau départemental.}

% ======================= RECHERCHE =======================
\section{Expérience de recherche}
\entree{Groupe de recherche KYMA (UNAL) — Chercheur, capteurs et commande}{Bogotá, 2024–présent}{%
KYMA est le groupe de recherche en capteurs, contrôle et intelligence artificielle du département de mécatronique de l'UNAL.
\begin{points}
  \item Projet \textit{Piezo} : synthèse et caractérisation du titanate de baryum (\bati) pour des capteurs piézoélectriques de force et de vibration sans plomb.
  \item Conception électronique pour l'acquisition des signaux et l'IoT ; algorithmes de commande robuste pour systèmes dynamiques.
  \item Gestion des ressources techniques du groupe.
\end{points}}

% ======================= PROJETS D'INGÉNIERIE =======================
\section{Projets d'ingénierie}

\entree{Aéropendule — identification et quatre lois de commande}{Técnicas de Control, 2026 · équipe de 4}{%
\begin{points}
  \item Identification expérimentale d'un pendule actionné par hélice (essais quasi statiques, oscillations libres, échelons) :
        modèle non linéaire validé sous Simulink avec \textbf{R\textsuperscript{2} = 0,943} et RMSE = 3,37°.
  \item Linéarisation à 70° et synthèse d'un PID par lieu des racines, puis d'un correcteur \textbf{H$\infty$} par sensibilité mixte (S, T, KS),
        d'une commande adaptative \textbf{MRAC} et d'une commande par \textbf{mode glissant}.
  \item Déploiement des quatre lois sur Arduino depuis Simulink et comparaison sur le banc réel (rejet de perturbation, robustesse au retard).
\end{points}
\outils{MATLAB/Simulink, Simulink Support Package for Arduino, C/C++.} \href{https://github.com/NicolasPlata2004/MATLAB_Pendulo}{\small[code]}}

\entree{Bras robotique plan 2R « ADANISS »}{Servomecanismos, 2026 · équipe de 5}{%
\begin{points}
  \item Trajectoire en trèfle paramétrée en polaire, cinématique inverse (coude haut), profils quintiques pour raccorder approche et cycle.
  \item Modèle lagrangien pour dimensionner les motoréducteurs (couples efficace et crête) ; transmission par courroie synchrone.
  \item Identification des articulations par essais en impulsion (\textbf{R\textsuperscript{2} > 0,99}) : gain, constante de temps, zone morte.
  \item PID discret sur ESP32, visualisation temps réel sous MATLAB à 60 Hz ; carte électronique (2 ponts BTS7960, capteurs de courant) conçue et fabriquée par l'équipe.
\end{points}
\outils{MATLAB, ESP32, Fusion 360, impression 3D PETG, découpe laser.} \href{https://github.com/NicolasPlata2004/ADANISS-2R}{\small[code]}}

\entree{Case Erector — formeuse de caisses automatisée}{Proyecto industrial, 2026 · équipe de 5}{%
\begin{points}
  \item Programmation en ladder (Studio 5000) d'une machine à trois axes à 20 caisses/min sur automate Allen-Bradley émulé.
  \item Validation sans machine réelle grâce à un jumeau numérique sous Siemens NX relié en OPC DA.
  \item Mesure du cycle dans les Trends (un cycle toutes les 3 s) et compensation de la latence de communication.
\end{points}
\outils{Studio 5000, automate émulé, OPC DA, Siemens NX.} \href{https://github.com/NicolasPlata2004/Case-Erector-Motion-control-of-a-three-axis-box-former-with-a-digital-twin}{\small[code]}}

\entree{ThermaLink-C3 — supervision de climatisation consciente de l'occupation}{Redes de Computadores, 2026 · individuel}{%
\begin{points}
  \item Objet embarqué sur ESP32-C3 : détection de présence (PIR), température (DHT22), clonage universel des trames infrarouges
        (protocoles courts et longs, validé sur trois climatiseurs de marques différentes avec un seul binaire).
  \item Supervision hiérarchique par machine à états et hystérésis ; firmware v2 avec panneau web (API REST), NTP, persistance NVS,
        \textbf{MQTT et découverte automatique Home Assistant}.
  \item Validation en boucle fermée contre un second nœud ESP32 émulant le climatiseur ; PCB deux couches sous KiCad, boîtier Fusion 360.
  \item Modèle énergétique : économie projetée de 25 à 30 \% de la climatisation, retour sur investissement inférieur à trois mois. Rédaction d'un article technique.
\end{points}
\outils{C++, ESP32-C3, MQTT, Home Assistant, KiCad, Fusion 360.}}

\entree{Passerelle IoT industrielle — commande d'entreprise NEXERGY}{IMT Atlantique, 2026–2027 · équipe de 7}{%
\begin{points}
  \item Responsable technique : déployer et qualifier la passerelle multiprotocole MultiEdge (MQTT, Modbus, OPC-UA, S7, EtherNet/IP) jusqu'au cloud.
  \item Banc d'essai de la chaîne compteur $\rightarrow$ MQTT $\rightarrow$ edge : tests de coupure du broker, reconnexion et perte de données.
\end{points}
\outils{Mosquitto, Python (paho-mqtt), automates Siemens et Allen-Bradley, OPC-UA.}}

\entree{Lévitateur magnétique analogique — rétro-ingénierie}{Projet personnel, 2026}{%
\begin{points}
  \item Reconstruction de la boucle de commande d'un lévitateur sans microcontrôleur : capteurs Hall différentiels, correcteur PD à avance de phase,
        étage de puissance push-pull ; analyse du principe (théorème d'Earnshaw).
  \item Série de neuf vidéos animées et bibliothèque Python pour dessiner des schémas électroniques.
\end{points}
\outils{Python, Manim.} \href{https://github.com/NicolasPlata2004/Magnetic-Levitator}{\small[code]}}

\entree{Contrôle thermique périopératoire}{Systèmes embarqués, UNAL}{%
\begin{points}
  \item Boucle PI sur microcontrôleur PIC18F4550 programmé en C ; traitement numérique du signal en temps réel.
\end{points}}

\entree{Autres projets techniques}{UNAL, 2021–2026}{%
\begin{points}
  \item \textbf{Dispensateur « Fluxia »} : script paramétrique Fusion 360 (API Python) générant l'assemblage complet (9 composants, 37 paramètres).
  \item \textbf{Capteur capacitif} : conception de PCB en équipe ; \textbf{système d'irrigation} en logique numérique (Verilog).
  \item \textbf{Modèles proie–prédateur} : Lotka–Volterra vs modèle écologique réaliste, simulation MATLAB/Python et linéarisation.
  \item \textbf{Solveur de systèmes linéaires} (Python) : Gauss–Jordan, Gauss–Seidel et Cramer avec traces pas à pas.
\end{points}}

% ======================= LOGICIEL =======================
\section{Projets logiciels}
\begin{points}
  \item \textbf{Eureka} — plateforme d'e-learning pour l'I.E. Luis María Jiménez (Aguazul) : numérisation d'examens par OCR local,
        apprentissage adaptatif et difficulté dynamique (React, Tesseract.js). Variante gamifiée pour l'algèbre de 8\textsuperscript{e} (Next.js).
  \item \textbf{Kairos} — application web progressive de suivi d'habitudes, synchronisée dans le cloud (React, Firebase).
  \item \textbf{François} — tuteur de français avec IA locale (Ollama), conversation et correction, 100 \% hors ligne.
  \item \textbf{Chaîne de transcription de cours} — faster-whisper (large-v3) sur GPU, puis génération de notes bilingues.
  \item \textbf{Bibliothèque Laplace pour TI-Nspire} — transformées et décomposition en éléments simples, fichier .tns généré depuis Python.
  \item \textbf{« Una gota más »} — application de gestion de banque de sang (Java), projet de structures de données dont j'ai coordonné l'équipe.
\end{points}

% ======================= ENSEIGNEMENT, VULGARISATION =======================
\section{Encadrement, prototypage et vulgarisation}
\entree{Moniteur académique}{UNAL}{Accompagnement d'étudiants et appui pédagogique.}
\entree{UN Makerspace — Prototypage et IoT}{Bogotá, 2023–présent}{%
Conception mécanique et prototypage par impression 3D ; automatisation de flux de travail (n8n) pour la traçabilité des projets.}
\entree{Vulgarisation scientifique}{2024–présent}{%
\begin{points}
  \item Séries d'animations pédagogiques programmées avec Manim : lévitateur magnétique (9 vidéos), aéropendule, équation du trèfle (vidéos en espagnol et en français).
  \item Fondateur du blog « El Éter » : articles techniques et analyses d'avancées technologiques.
  \item Vidéos de présentation de projets pour LinkedIn (lévitateur, Case Erector, ThermaLink-C3).
\end{points}}

% ======================= COMPÉTENCES =======================
\section{Compétences techniques}
{\small
\begin{tabularx}{\linewidth}{@{}>{\bfseries}p{3.6cm}X@{}}
Automatique & Identification, représentation d'état, lieu des racines, Bode/Nyquist, PID, H$\infty$, MRAC, mode glissant, commande numérique. \\[2pt]
Simulation & MATLAB, Simulink (déploiement sur Arduino), Python (NumPy, SciPy, Matplotlib), LTspice, Proteus, Multisim. \\[2pt]
Embarqué \& IoT & ESP32 / ESP32-C3, Arduino, PIC18F4550, C/C++, MQTT, Home Assistant, API REST, capteurs, KiCad (PCB). \\[2pt]
Automatisme & Studio 5000 (ladder), automates émulés, OPC DA/UA, Modbus, jumeau numérique (Siemens NX), CADe SIMU, FluidSIM. \\[2pt]
CAO / FAO & Fusion 360 (dont API Python), Inventor, AutoCAD, Mastercam, impression 3D, découpe laser. \\[2pt]
Programmation & Python (Pandas, NumPy), MATLAB, C/C++, Java, JavaScript/TypeScript, React, Next.js, Firebase, Git, LaTeX. \\[2pt]
Communication & Manim, montage vidéo, rédaction technique en espagnol, français et anglais. \\
\end{tabularx}}

% ======================= LANGUES =======================
\section{Langues}
{\small
\begin{tabularx}{\linewidth}{@{}>{\bfseries}p{3.6cm}X@{}}
Espagnol & Langue maternelle ; excellente rédaction technique. \\
Français & B2 — certification officielle ; études d'ingénieur en français. \\
Anglais & C1 (certification B2) ; communication technique courante. \\
Allemand & A1 — notions. \\
\end{tabularx}}

% ======================= CERTIFICATIONS =======================
\section{Cours et certifications}
\begin{points}
  \item Diplôme en développement d'applications mobiles — MinTIC / Universidad Tecnológica de Pereira.
  \item Fondamentaux de la programmation et du développement logiciel — MinTIC.
  \item Introduction à l'impression 3D métallique.
  \item Introduction à Python pour l'analyse de données.
\end{points}

% ======================= DIVERS =======================
\section{Centres d'intérêt et références}
{\small
\textbf{Intérêts :} vulgarisation scientifique, impression 3D et électronique, apprentissage des langues.\\
\textbf{Références :} lettre de recommandation du Pr José Manuel Arroyo Osorio (Faculté d'ingénierie, UNAL) ; coordonnées sur demande.}

\end{document}
